\section{A near-threshold obstruction to a pointwise interaction-virial sign}
\label{sec:r9-virial-obstruction}

This section concerns the pure mass-critical BNLS
\[
 iu_t=\Delta^2u-|u|^{8/N}u,\qquad N\ge1.
\]
Write
\[
 p=2+8/N,\quad M(f)=\|f\|_2^2,\quad A(f)=\|\Delta f\|_2^2,
 \quad E(f)=\frac12A(f)-\frac1p\|f\|_p^p,
\]
\[
 P(f)=\operatorname{Im}\int\overline f\nabla f,
 \qquad W(f)=\operatorname{Im}\int\Delta f\nabla\overline f.
\]
Let $M_c$ be the sharp Gagliardo--Nirenberg critical mass. Fix the
interaction cutoff $\phi_R$ which equals $|z|^2/2$ for $|z|\le4R$,
is constant for $|z|\ge8R$, and satisfies $D^2\phi_R\le I$. Define
\begin{equation}\label{eq:r9-interaction-functional}
 I_R(f)=2\iint |f(y)|^2\nabla\phi_R(x-y)\cdot
       \operatorname{Im}(\overline f(x)\nabla f(x))\,dx\,dy.
\end{equation}
For smooth compactly supported $f$, let $\mathcal D_R(f)$ denote
the initial derivative of this functional prescribed by the BNLS
vector field, equivalently the right-hand side of the exact
interaction-virial identity evaluated at $f$.

\begin{theorem}[Positive localized virial derivative arbitrarily near threshold]
\label{thm:r9-static-obstruction}
For every $\delta>0$, there exist constants
$M_*,A_*,e_*,c_*,R_*>0$ and functions
$f_R\in C_c^\infty(\mathbb R^N;\mathbb C)$, $R\ge R_*$, such that
\begin{align}
 M_c<M(f_R)&=M_*<M_c+\delta,&
 A(f_R)&=A_*,\label{eq:r9-static-invariants}\\
 P(f_R)&=0,&E(f_R)&=-e_*<0,\notag
\end{align}
and
\begin{equation}\label{eq:r9-static-positive}
 \mathcal D_R(f_R)\ge c_*>0.
\end{equation}
All constants are independent of $R$. Moreover, one may choose
these families as $\delta\downarrow0$ so that, for a fixed real
normalized optimizer $Q$,
\[
 \sup_{R\ge R_*(\delta)}\|f_R-Q\|_{H^2}\longrightarrow0.
\]
Consequently no estimate
\[
 \mathcal D_R(f)\le c\,M(f)E(f)+o_{R\to\infty}(1),\qquad c>0,
\]
can hold uniformly over zero-momentum negative-energy data of
sufficiently near-critical mass and bounded kinetic energy.
\end{theorem}

\begin{lemma}[Independent momentum and flux directions]
\label{lem:r9-independent-currents}
For every nonzero real $q\in C_c^\infty(\mathbb R^N)$, the linear
map
\[
 b\longmapsto \bigl(P(q+ib),W(q+ib)\bigr),
 \qquad b\in C_c^\infty(\mathbb R^N;\mathbb R),
\]
is onto $\mathbb R^N\times\mathbb R^N$.
\end{lemma}
\begin{proof}
Integration by parts gives
\begin{equation}\label{eq:r9-current-linearization}
 P(q+ib)=-2\int b\nabla q,\qquad
 W(q+ib)=2\int b\nabla\Delta q.
\end{equation}
The $2N$ real functions $\partial_jq,\partial_j\Delta q$ are
linearly independent in $L^2$. Indeed, a dependence would imply
\[
 \bigl[(a\cdot\xi)-(b_0\cdot\xi)|\xi|^2\bigr]\widehat q(\xi)=0
 \quad\text{a.e.}
\]
for some $a,b_0\in\mathbb R^N$. A nonzero polynomial has a
Lebesgue-null zero set, whereas a nonzero $L^2$ function cannot
be supported on a null set. Thus the polynomial vanishes
identically, forcing $a=b_0=0$. The Gram matrix of the $2N$
functions is therefore invertible. Taking $b$ in their span
realizes every pair in \eqref{eq:r9-current-linearization}.
\end{proof}

\begin{lemma}[Separated-packet identity]
\label{lem:r9-separated-virial}
Suppose $F,G\in C_c^\infty$ have support diameters below $4R$, and
the distance between their supports is above $8R$. Then
\begin{equation}\label{eq:r9-separated-identity}
 \mathcal D_R(F+G)=\mathcal V(F)+\mathcal V(G),\qquad
 \mathcal V(h)=16M(h)E(h)-8P(h)\cdot W(h).
\end{equation}
\end{lemma}
\begin{proof}
For intrapacket pairs the cutoff is exactly quadratic; for
interpacket pairs it is constant. Substitution into the exact
interaction identity gives \eqref{eq:r9-separated-identity}.
Equivalently, the quadratic interaction associated with one
packet is
\[
 2\left[M(h)\operatorname{Im}\int x\cdot\overline h\nabla h
       -\left(\int x|h|^2\right)\cdot P(h)\right].
\]
With $J(h)=\operatorname{Im}\int x\cdot\overline h\nabla h$,
one has $J'=8E(h)$, while the mass first moment has derivative
$4W(h)$. Mass and momentum are conserved, so the displayed
interaction has derivative $2[8ME-4W\cdot P]$, as asserted. This is an
initial-derivative calculation; no finite-speed propagation or
persistence of compact support is asserted.
\end{proof}

\begin{proof}[Proof of Theorem~\ref{thm:r9-static-obstruction}]
Fix a real optimizer $Q$ with $M(Q)=M_c$ and $E(Q)=0$.
Approximate it in $H^2$ by a nonzero real compactly supported
smooth function $q_0$. Multiplication by
\[
 \left(\frac{pA(q_0)}{2\|q_0\|_p^p}\right)^{1/(p-2)}
\]
gives a real $q\in C_c^\infty$ satisfying
\begin{equation}\label{eq:r9-zero-energy-core}
 E(q)=0,\qquad M_0:=M(q)<M_c+\delta/4.
\end{equation}
The multiplier tends to one as $q_0\to Q$.

By Lemma~\ref{lem:r9-independent-currents}, there are fixed real
$b_1,b_2\in C_c^\infty$ such that
\[
 P(q+ib_1)=0,\quad W(q+ib_1)=e_1,\qquad
 P(q+ib_2)=-e_1,\quad W(q+ib_2)=0.
\]
Set $b_\varepsilon=\varepsilon b_1+\varepsilon^2b_2$. Then
\begin{equation}\label{eq:r9-prescribed-core-currents}
 P(q+ib_\varepsilon)=-\varepsilon^2e_1,\qquad
 W(q+ib_\varepsilon)=\varepsilon e_1,\qquad
 \|b_\varepsilon\|_{H^2}=O(\varepsilon).
\end{equation}
Since
\[
 \left.\partial_aE(aq)\right|_{a=1}
 =A(q)-\|q\|_p^p=-\frac4N A(q)<0,
\]
the implicit-function theorem supplies $a_\varepsilon=1+O(\varepsilon^2)$
for which
\begin{equation}\label{eq:r9-core-negative-energy}
 F_\varepsilon=a_\varepsilon q+ib_\varepsilon,
 \qquad E(F_\varepsilon)=-\frac{\varepsilon^3}{4M_0}.
\end{equation}
Here the energy is $C^2$ on this finite-dimensional family because
$p>2$, and its change under the imaginary perturbation starts at
second order. Write $M_F=M(F_\varepsilon)$. Exactly,
\[
 M_F=a_\varepsilon^2M_0+\|b_\varepsilon\|_2^2,
 \qquad P_F=-a_\varepsilon\varepsilon^2e_1,
 \qquad W_F=a_\varepsilon\varepsilon e_1.
\]
Consequently
\begin{equation}\label{eq:r9-core-positive-virial}
 \mathcal V(F_\varepsilon)
 =\left(8a_\varepsilon^2-\frac{4M_F}{M_0}\right)\varepsilon^3
 =(4+O(\varepsilon^2))\varepsilon^3>2\varepsilon^3
\end{equation}
for all sufficiently small positive $\varepsilon$.

Fix $\mathfrak m\in(0,\min\{M_c,\delta/4\})$, and choose a real
radial $\chi\in C_c^\infty(\mathbb R^N)$ with $\|\chi\|_2=1$.
Define a second packet by
\begin{equation}\label{eq:r9-balancing-packet}
 G(x)=\sqrt{\mathfrak m}\,L^{-N/2}\chi(x/L)e^{ikx_1},
 \qquad k=\frac{a_\varepsilon\varepsilon^2}{\mathfrak m}.
\end{equation}
Its mass is $\mathfrak m$ and its momentum is exactly $-P_F$.
Put
\[
 B=\|\nabla\chi\|_2^2,\quad C=\|\Delta\chi\|_2^2,
 \quad D=\|\chi\|_p^p,\quad d=1+2/N.
\]
Fourier expansion and radial symmetry give
\begin{align}
 A(G)&=\mathfrak m\bigl[k^4+2d k^2L^{-2}B+L^{-4}C\bigr],
 \label{eq:r9-radiation-formulas}\\
 W(G)&=\mathfrak m\bigl[k^3+d kL^{-2}B\bigr]e_1,
 &\|G\|_p^p&=\mathfrak m^{p/2}L^{-4}D.\notag
\end{align}
In particular, the plane-wave terms cancel exactly in the
self-interaction derivative:
\begin{equation}\label{eq:r9-radiation-positive-virial}
 \mathcal V(G)=8\mathfrak m^2
 \left[d k^2L^{-2}B+
 L^{-4}\left(C-\frac2p\mathfrak m^{4/N}D\right)\right]>0.
\end{equation}
The last inequality follows from the sharp GN inequality and
$\mathfrak m<M_c$.
For fixed $\varepsilon$, as $L\to\infty$,
\[
 E(G)\longrightarrow\frac{\mathfrak m k^4}{2}
 =\frac{a_\varepsilon^4\varepsilon^8}{2\mathfrak m^3}.
\]
Choose $\varepsilon>0$ small enough that this limit is less than
$\varepsilon^3/(16M_0)$ and that $M_F<M_c+\delta/2$.
Then choose $L$ sufficiently large that
\begin{equation}\label{eq:r9-radiation-energy-budget}
 0<E(G)<\frac{\varepsilon^3}{8M_0}.
\end{equation}

Fix these packets. For sufficiently large $R$, set
\[
 f_R(x)=F_\varepsilon(x)+G(x-20Re_1).
\]
Their support diameters are below $4R$ and their mutual support
distance exceeds $8R$. All invariants in
\eqref{eq:r9-static-invariants} are exactly independent of $R$
by disjointness of the supports and locality of the pure BNLS
energy. Their momentum is zero, and
\[
 E(f_R)=E(F_\varepsilon)+E(G)<-\frac{\varepsilon^3}{8M_0}<0,
 \qquad M(f_R)<M_c+\delta.
\]
Negative energy forces $M(f_R)>M_c$. Lemma~\ref{lem:r9-separated-virial}
and \eqref{eq:r9-core-positive-virial}--\eqref{eq:r9-radiation-positive-virial}
prove \eqref{eq:r9-static-positive} with $c_*=2\varepsilon^3$.
Finally, as $\delta\downarrow0$, choose $q\to Q$,
$\mathfrak m\to0$, and then successively smaller $\varepsilon$
and larger $L$. The core tends to $Q$ in $H^2$, and
\eqref{eq:r9-radiation-formulas} makes $\|G\|_{H^2}\to0$.
This norm is translation invariant, proving uniform closeness
in $R$.
\end{proof}

\begin{remark}[The exact signed remainder]
\label{rem:r9-hessian-deficit}
The obstruction persists when the full favorable Hessian deficit
is retained. Let $\mathsf D_R=I-D^2\phi_R\ge0$, and set
$\rho=|f|^2$, $j=\operatorname{Im}(\overline f\nabla f)$,
$w=\operatorname{Im}(\Delta f\nabla\overline f)$. Define
\begin{align*}
 \mathcal K_R&=\iint\rho(y)(\mathsf D_R)_{jk}(x-y)
  \operatorname{Re}(\partial_{ij}\overline f\,\partial_{ik}f)(x)
  \,dx\,dy,\\
 \mathcal N_R&=\frac{2\beta}{p}\iint\rho(y)
  \operatorname{tr}\mathsf D_R(x-y)|f(x)|^p\,dx\,dy,
  \qquad\beta=8/N,\\
 \mathcal X_R&=8\iint\mathsf D_R(x-y)[j(x),w(y)]\,dx\,dy.
\end{align*}
For the support configurations above, every term containing a
higher cutoff derivative vanishes exactly, and
\[
 \mathcal D_R(f_R)=16ME-8P\cdot W
                  -8\mathcal K_R+\mathcal N_R+\mathcal X_R.
\]
Moreover,
\begin{align*}
 \mathcal K_R&=M_F A(G)+\mathfrak m A(F_\varepsilon),\\
 \mathcal N_R&=\frac{16}{p}
  \bigl(M_F\|G\|_p^p+\mathfrak m\|F_\varepsilon\|_p^p\bigr),\\
 \mathcal X_R&=8\bigl(P_F\cdot W(G)+P(G)\cdot W_F\bigr).
\end{align*}
Since $P=0$ and $\mathcal D_R(f_R)>0>16ME$, one has the strict
inequality $\mathcal N_R+\mathcal X_R>8\mathcal K_R$. Thus the
failure of a universal pointwise sign is an exact property of
the interaction identity. The theorem imposes no conclusion on
the subsequent lifespan or long-time dynamics of these data.
\end{remark}
